\section{Discussion and Conclusion}
\label{sec:conclusion}

Zero-shot cross-embodiment transfer for tabletop manipulation with two-finger grippers decomposes into four channels through which the source robots leave their mark on a policy, and the controlled study measures each of them under one backbone, one budget and one evaluation. The protocol channel comes first: the target-absent and target-seen settings differ by up to eighteen points for the same representation, and only the target-absent setting reveals the full effect of the other channels. On the representation channel the frame of the action orders the outcomes, from the base frame that stores a convention, through the tool frame that stores a geometry, to the camera frame that stores an outcome; the camera-frame effect transfers best, and its prediction error on a new robot forecasts how well that robot will do. On the data channel a diverse source family under a fixed budget substitutes for source-specific structure, and effect supervision at eight sources matches tool-frame commands at 128. On the supervision channel image-grounded targets help most, and the effect target leads them. These findings hold for fixed-base arms with two-finger grippers in tabletop manipulation with matched cameras, which is the regime the benchmark was built to isolate. Grippers with more fingers, mobile bases and long-horizon tasks change the set of keypoints, the camera alignment and the horizon of the effect, and extending the study to them is the natural next step. The measured results on established benchmarks show that the recipe keeps its value beyond the controlled setting, and the level result on a physically different arm with proxy anchors and unmatched cameras shows that the gain scales with how closely the cameras and anchors of the deployment match those of training. We therefore recommend four practices for cross-embodiment evaluation: report whether the target was seen during pretraining, express the action in a frame the target shares and prefer the camera frame when cameras can be matched, treat source diversity and budget as coupled variables, and measure the effect error on the target before spending task rollouts on it.

\subsection*{AI use statement}
Generative AI tools assisted with code refactoring, log analysis, literature search and language editing of this manuscript. All experimental designs, protocols, analyses and claims were specified and verified by the authors. AI-generated code was reviewed and tested by the authors before use. We take responsibility for the final content of this work.

\subsection*{Ethics statement}
This work uses simulated robots and publicly released robot demonstration datasets, and all data are free of personal information. The methods aim at reusable manipulation policies and carry the general risks of autonomous robot control, which we mitigate through evaluation in simulation before deployment.

\subsection*{Reproducibility statement}
Sections~\ref{sec:method} and \ref{sec:design} specify the protocols, the representations, the objective and the decoder; Section~\ref{sec:bench} specifies the benchmark; Appendix~\ref{app:family} lists the generator of the source family, Appendix~\ref{app:hparams} the training and evaluation hyperparameters, Appendix~\ref{app:prereg} the pre-registered design and Appendix~\ref{app:compute} the compute budget. The embodiment generator, the benchmark scenes, the training configurations and the evaluation scripts are released with the paper.
